\documentclass[11pt]{amsart}
\usepackage{geometry}                % See geometry.pdf to learn the layout options. There are lots.
\geometry{letterpaper}                   % ... or a4paper or a5paper or ... 
%\geometry{landscape}                % Activate for for rotated page geometry
%\usepackage[parfill]{parskip}    % Activate to begin paragraphs with an empty line rather than an indent
\usepackage{graphicx}
\usepackage{amssymb}
\usepackage{epstopdf}
\DeclareGraphicsRule{.tif}{png}{.png}{`convert #1 `dirname #1`/`basename #1 .tif`.png}

\title{PS 2.12}
\author{Enrique Zuñiga Zuñiga}
%\date{}                                           % Activate to display a given date or no date

\begin{document}
\maketitle
%\section{}
\noindent Construya un diagrama de flujo tal que dado como dato una temperatura en grados Fahrenheit, determine el deporte que es apropiado practicar a esa temperatura, teniendo en cuenta la siguiente tabla:\newline
$$
\begin{array}{|c|c|}
\hline
Tabla 6.1 &\\
\hline
DEPORTE & TEMPERATURA\\
\hline
Natacion & > 85\\
\hline
Tenis &70 < TEMP <= 85\\
\hline
Golf &32 < TEMP <= 70\\
\hline
Esqui &10 < TEMP <= 32\\
\hline
Marcha &<= 10\\
\hline
\end{array}
$$\newline

\noindent Dato: TEMP (variable de tipo real que representa la temperatura).\newline


\end{document}  